% Figure from MATH 590 notes, section 26. Compile with the notes preamble.
\begin{tikzpicture}[scale=1.2, x={(1cm,0cm)}, y={(0.55cm,0.38cm)}, z={(0cm,1cm)}]
  % z-axis below the plane (behind)
  \draw[figR, thick, dashed] (0,0,-1.6) -- (0,0,0);
  % the plane z = 0 minus the origin
  \filldraw[fill=figB!14, draw=figB, thick] (-2,-1.8,0) -- (2,-1.8,0) -- (2,1.8,0) -- (-2,1.8,0) -- cycle;
  \node[font=\scriptsize, figB, anchor=north west] at (-2,-1.8,0) {plane $z = 0$, minus $0$};
  % vertical arrows (x,y,z) -> (x,y,0)
  \foreach \p/\q/\h in {1.3/0.9/1.3, -1.2/1.1/1.0, -0.9/-1.0/1.45, 0.9/-0.9/0.8} {
        \draw[figR!85, thick, -{Stealth[length=4pt]}] (\p,\q,\h) -- (\p,\q,0.1);
    \fill (\p,\q,\h) circle (1pt);
    \fill[figB] (\p,\q,0) circle (1pt);
  }
  \node[font=\scriptsize, anchor=west] at (1.3,0.9,1.3) {$(x,y,z)$};
  \node[font=\scriptsize, figB, anchor=north east] at (1.36,0.9,-0.06) {$(x,y,0)$};
  % z-axis above the plane (front), removed
  \draw[figR, thick, dashed] (0,0,0) -- (0,0,1.9) node[above, font=\scriptsize] {$z$-axis removed};
  \filldraw[fill=white, draw=figR, thick] (0,0,0) circle (1.6pt);
\end{tikzpicture}
